% EconPaper Architect · LaTeX 论文模板
% 编译：xelatex -> bibtex -> xelatex -> xelatex
% 说明：中文用 ctex（XeLaTeX）；英文投稿时删掉 \usepackage{ctex} 并改 documentclass 为对应 cls。
\documentclass[11pt,a4paper]{article}

% ---------- 中文与字体 ----------
\usepackage{ctex}
\usepackage{geometry}
\geometry{left=3cm,right=3cm,top=2.5cm,bottom=2.5cm}

% ---------- 数学 ----------
\usepackage{amsmath,amssymb,amsthm,mathtools}
\numberwithin{equation}{section}
\theoremstyle{plain}
\newtheorem{hyp}{假说}

% ---------- 表格：一律三线表 ----------
\usepackage{booktabs}
\usepackage{threeparttable}
\usepackage{multirow}
\usepackage{longtable}
\usepackage{siunitx}
\sisetup{input-symbols=(), detect-all}

% ---------- 图形 ----------
\usepackage{graphicx}
\usepackage{caption}
\captionsetup{font=small, labelfont=bf}
\graphicspath{{figures/}}

% ---------- 引用 ----------
\usepackage[round]{natbib}
\bibliographystyle{aer}          % 中文期刊可换 plainnat / gbt7714-numerical
\usepackage{hyperref}
\hypersetup{colorlinks=true,citecolor=blue,linkcolor=black,urlcolor=blue}

% ---------- 版式 ----------
\usepackage{setspace}
\onehalfspacing
\usepackage{indentfirst}
\setlength{\parskip}{0.3em}
\usepackage{titlesec}
\titleformat{\section}{\large\bfseries}{\thesection}{1em}{}

% ---------- 自定义命令 ----------
\newcommand{\sym}[1]{\rlap{$^{#1}$}}        % 显著性标注：\sym{***}
\newcommand{\cneed}[1]{\textbf{[待填：#1]}}

% ============================================================
\title{[论文标题：中文 $\le$ 25 字 / 英文 $\le$ 15 词]}
\author{[作者]\thanks{[基金资助 / 通讯作者 / 致谢。AI 辅助情况须按目标期刊要求披露。]}}
\date{\today}

\begin{document}
\maketitle

% ---------------- 摘要 ----------------
\begin{abstract}
\noindent [200--300 字（英文 150--250 词）。五要素：问题 $\rightarrow$ 方法 $\rightarrow$ 数据 $\rightarrow$ 主要发现（含量级）$\rightarrow$ 贡献。不引用文献、不含公式、不含脚注。]

\par\smallskip
\noindent\textbf{关键词：} [关键词1；关键词2；关键词3]

\par\smallskip
\noindent\textbf{JEL 分类号：} [如 J23, C23]
\end{abstract}

% ---------------- 正文 ----------------
\section{引言}
[1200--1800 字。AER 式七段式：钩子 $\rightarrow$ 研究问题 $\rightarrow$ 识别策略 $\rightarrow$ 主要发现（含量级）$\rightarrow$ 文献贡献 $\rightarrow$ 政策相关性 $\rightarrow$ 路线图。]

\section{文献综述}
[中文核心式独立成节；AER 式并入引言。按主题分组，每组末尾指出本文边际贡献。]

\section{理论框架}
\subsection{模型设定}
[假设呈现顺序：从最一般到最限制性。]

\begin{equation}\label{eq:model}
\max_{c_t}\ \sum_{t=0}^{\infty}\beta^{t}u(c_t)
\quad \text{s.t.} \quad
a_{t+1}=(1+r)a_t+w_t-c_t
\end{equation}

\subsection{假说推导}
\begin{hyp}\label{hyp:h1}
由式~\eqref{eq:model} 可得：[假说内容，须与实证部分变量形成可检验对应关系]。
\end{hyp}

\section{数据与变量}
\subsection{描述性统计}
% 三线表示例（由 esttab / fsum 导出）
\begin{table}[htbp]
\centering
\begin{threeparttable}
\caption{描述性统计}
\label{tab:desc}
\begin{tabular}{lccccc}
\toprule
变量 & 均值 & 标准差 & 最小值 & 最大值 & 观测值 \\
\midrule
被解释变量 $Y$ & 0.150 & 0.120 & 0.010 & 0.890 & 3,240 \\
核心解释变量 $X$ & 0.450 & 0.230 & 0.000 & 1.000 & 3,240 \\
\bottomrule
\end{tabular}
\begin{tablenotes}
\footnotesize
\item 注：数据来源为\cneed{填写来源}；样本区间为\cneed{填写区间}。存在 \texttt{[SIM]} 标记者为模拟数据，仅供格式演示。
\end{tablenotes}
\end{threeparttable}
\end{table}

\subsection{制度背景}
[政策时点、实施范围、执行强度差异——DID / RDD 识别成立的事实基础。]

\section{实证策略}
\subsection{基准模型}
\begin{equation}\label{eq:did}
Y_{it}=\alpha+\beta D_{it}+\gamma X_{it}+\mu_i+\lambda_t+\varepsilon_{it}
\end{equation}
其中 $D_{it}$ 为处理变量，$\mu_i$ 与 $\lambda_t$ 分别为个体与时间固定效应，标准误聚类到\cneed{层级}层面。

\subsection{识别假设}
[硬阻断项 H3：本节缺失则校验不通过。须陈述假设、论证其在本文情境下成立、并给出证伪检验。]

% 回归结果三线表示例
\begin{table}[htbp]
\centering
\begin{threeparttable}
\caption{基准回归结果}
\label{tab:main}
\begin{tabular}{lccc}
\toprule
 & (1) & (2) & (3) \\
\midrule
核心解释变量 & 0.150\sym{***} & 0.132\sym{***} & 0.128\sym{***} \\
 & (0.021) & (0.020) & (0.019) \\
\midrule
控制变量 & 否 & 是 & 是 \\
个体固定效应 & 是 & 是 & 是 \\
时间固定效应 & 是 & 是 & 是 \\
\midrule
样本量 & 3,240 & 3,240 & 3,240 \\
调整 $R^2$ & 0.112 & 0.203 & 0.211 \\
\bottomrule
\end{tabular}
\begin{tablenotes}
\footnotesize
\item 注：括号内为聚类到个体层面的稳健标准误；\sym{*} $p<0.1$，\sym{**} $p<0.05$，\sym{***} $p<0.01$。
\end{tablenotes}
\end{threeparttable}
\end{table}

\section{结果}
[用具体数字陈述：``$X$ 每增加 1 个标准差，$Y$ 增加 0.15 个单位（$p<0.01$）''。]

\begin{figure}[htbp]
\centering
\includegraphics[width=0.8\textwidth]{fig_event.png}
\caption{事件研究图\quad{\footnotesize 注：数据来源于\cneed{填写}；阴影为 95\% 置信区间。}}
\label{fig:event}
\end{figure}

\section{稳健性检验}
[由 \texttt{scripts/decide.py} 输出的稳健性清单填入。]

\section{结论与政策建议}
[主要发现 $\rightarrow$ 政策含义（基于结果，不凭空得出）$\rightarrow$ 研究局限 $\rightarrow$ 未来方向。]

% ---------------- 参考文献 ----------------
\bibliography{references}

% ---------------- 附录 ----------------
\appendix
\section{证明}
\section{额外表格}
\section{数据说明}
\section{代码}

\end{document}
